\appendix

\section{More about operators on a Hilbert space}

We group in this appendix a series of definitions and theorems that are not necessarily part
of the basic foundations of quantum mechanics, but are still crucial in understanding how
exactly some computations are carried out and why they hold true. Some of the results
presented here go beyond the scope of our work as far as their proof is concerned, and only
serve the purpose of giving a clearer picture of what some tools actually are without having
to reference an entire course on functional analysis.

\textcolor{red}{TODO:}
\begin{itemize}
	\item introduce UNITARIES and ISOMETRIES (and talk more about self-adjoint operators?)
	\item ORTHOGONAL COMPLEMENT theorem + RIESZ representation theorem
	\item SPECTRAL theorem (which version?)
	\item DIRECT SUMS and TENSOR PRODUCTS of hilbert spaces (and their uniqueness)\\
		$\longrightarrow$ move N-fold tensor product state space definition/explanation here,
		and move construction of (anti-)symmetric tensor products into the canonical/grand
		canonical discussion, it makes more sense.
	\item INTEGRALS (and derivatives) of operator-valued functions
	\item comments on and PROPERTIES of $e^{-\beta H_0}$ (with $H_0$ bounded and maybe even
		unbounded)
	\item a couple of comments on VARIATIONAL PROBLEMS? (from blanchard maybe)
\end{itemize}

\section{On particle-hole excitations}

\textcolor{red}{How do we apply all this theory to, for example, BCS superconductors - since
Bogolubov transformations do not necessarily preserve the relations
$[a_\pm,a_\pm]_\mp = 0$ and $[a^*_\pm,a^*_\pm]_\mp = 0$ and we did rely on those too at times.}

\section{References}
\cite{test}
